\section{HIPÓTESIS}

\subsection{Planteamiento de la Hipótesis}
Considerando el problema principal del actual trabajo de investigación se plantea la siguiente hipótesis.

H1.``El protocolo de votación electrónica basado en criptografía post-cuántica de retículos coadyuvará a mantener la integridad y la confidencialidad del proceso de votación, verificadas formalmente mediante Tamarin Prover, con una latencia que no exceda en más del 30\% a la de un protocolo equivalente basado en RSA/ECDSA.''



\subsection{Definición de Variables}
De acuerdo a la hipótesis podemos sacar las siguientes variables:

\begin{itemize}
	\item \textbf{Variable independiente:} El protocolo de votación electrónica basado en criptografía post-cuántica de retículos.
	\item \textbf{Variable dependiente 1:} Mantenimiento de la integridad y la confidencialidad del proceso de votación.
	\item \textbf{Variable dependiente 2:} Latencia del protocolo, con un umbral de aceptación de hasta el 30\% adicional respecto al protocolo basado en RSA/ECDSA.
\end{itemize}

\newpage
\subsection{Matriz de Operacionalización de Variables}
\begingroup
\renewcommand{\arraystretch}{1.3}
\scriptsize
\setlength{\tabcolsep}{3pt}
\begin{longtable}{|>{\raggedright\arraybackslash}p{2.0cm}|>{\raggedright\arraybackslash}p{1.6cm}|>{\raggedright\arraybackslash}p{2.2cm}|>{\raggedright\arraybackslash}p{2.4cm}|>{\raggedright\arraybackslash}p{1.8cm}|>{\raggedright\arraybackslash}p{1.9cm}|>{\raggedright\arraybackslash}p{1.9cm}|}
	\caption{Matriz de operacionalización de variables.}
	\label{tab:operacionalizacion}\\
	\hline
	\textbf{Variable} &
	\textbf{Tipo} &
	\textbf{Definición Conceptual} &
	\textbf{Definición Operacional} &
	\textbf{Indicadores} &
	\textbf{Unidad / Valor} &
	\textbf{Instrumento / Fuente} \\
	\hline
	El protocolo de votación electrónica basado en criptografía post-cuántica de retículos &
	Independiente &
	Esquema criptográfico sobre el cual se sustenta el protocolo de votación electrónica &
	Naturaleza del esquema criptográfico aplicado al protocolo diseñado &
	Tipo de protocolo empleado &
	Categórica: post-cuántico (FIPS 203/204) / clásico (RSA/ECDSA) &
	Diseño experimental comparativo / implementación del protocolo \\
	\hline
	Mantenimiento de la integridad y la confidencialidad del proceso de votación &
	Dependiente &
	Correspondencia entre el voto emitido y el contabilizado, y secreto del contenido del voto frente a partes no autorizadas &
	Cumplimiento formal de las propiedades de integridad y confidencialidad en el modelo del protocolo &
	Cumplimiento de los lemmas de integridad y de confidencialidad &
	Verdadero / Falso &
	Tamarin Prover \\
	\hline
	Latencia del protocolo &
	Dependiente &
	Comportamiento del protocolo en el tiempo requerido para ejecutar sus operaciones criptográficas &
	Medición empírica de la latencia; se considera aceptable cuando no exceda en más del 30\% a la del protocolo RSA/ECDSA &
	Tiempo de firma, verificación y encapsulamiento &
	ms (umbral: $\le$ 30\% adicional respecto a RSA/ECDSA) &
	Implementación y simulación con distinto número de votantes \\
	\hline
\end{longtable}
\par\vspace{-12pt}
\begin{center}
	\footnotesize\textit{Fuente: Elaboración propia.}
\end{center}
\endgroup